\makeatletter
\@ifundefined{keepfig}{\newenvironment{keepfig}{\par\edef\keepfig@pi{\the\parindent}\noindent
  \begin{minipage}{\linewidth}\parindent=\keepfig@pi\relax}{\end{minipage}\par}}{}
\makeatother
\section{Теорема о примитивном элементе}

Напомним, что конечное расширение полей $E/F$ называется \emph{простым}, если $E=F[\alpha]$ для некоторого $\alpha\in E$;
такой элемент $\alpha$ называется \emph{примитивным элементом} расширения. Мы покажем, что все \emph{конечные}
сепарабельные расширения (и даже чуть больше) просты.

\begin{example}
Рассмотрим $E=\Q[\sqrt2,\sqrt3]$. Это поле разложения многочлена $(X^2-2)(X^2-3)$, и $[E:\Q]=4$. Действительно,
$\sqrt3\notin\Q[\sqrt2]$: из $\sqrt3=a+b\sqrt2$ ($a,b\in\Q$) следует $3=a^2+2b^2+2ab\sqrt2$, т.\,е.\ $ab=0$ (так как
$\sqrt2\notin\Q$) и $a^2=3$ или $2b^2=3$, что невозможно в $\Q$; поэтому $[E:\Q]=[E:\Q[\sqrt2]]\cdot[\Q[\sqrt2]:\Q]=2\cdot2$
(ср.\ упражнение 1-2). Значит, $E/\Q$~--- расширение Галуа с группой порядка $4$. Каждый автоморфизм переводит $\sqrt2\mapsto\pm\sqrt2$,
$\sqrt3\mapsto\pm\sqrt3$, а таких наборов знаков ровно четыре; значит, $\Gal(E/\Q)=\{1,\sigma,\tau,\sigma\tau\}$~---
четверная группа, где
\[
\begin{cases}\sigma\sqrt2=-\sqrt2,\\ \sigma\sqrt3=\sqrt3,\end{cases}\qquad
\begin{cases}\tau\sqrt2=\sqrt2,\\ \tau\sqrt3=-\sqrt3.\end{cases}
\]
Тогда
\[
\sigma(\sqrt2+\sqrt3)=-\sqrt2+\sqrt3,\qquad \tau(\sqrt2+\sqrt3)=\sqrt2-\sqrt3,\qquad
\sigma\tau(\sqrt2+\sqrt3)=-\sqrt2-\sqrt3,
\]
и все три числа отличны от $\sqrt2+\sqrt3$. Значит, подгруппа $\Gal(E/\Q[\sqrt2+\sqrt3])$, состоящая из автоморфизмов,
оставляющих $\sqrt2+\sqrt3$ на месте, тривиальна, и по основной теореме (\rf{3.17}) $\Q[\sqrt2+\sqrt3]=E^{\{1\}}=E$:
\[
\Q[\sqrt2,\sqrt3]=\Q[\sqrt2+\sqrt3].
\]
Попутно получаем минимальный многочлен примитивного элемента: это произведение множителей $X-\gamma'$ по всем
сопряжённым $\gamma'$, т.\,е.\ $\prod_{\pm,\pm}\bigl(X-(\pm\sqrt2\pm\sqrt3)\bigr)=X^4-10X^2+1$.
\end{example}

Ясно, что это рассуждение должно работать гораздо шире. Напомним, что алгебраический над $F$ элемент
\emph{сепарабелен} над $F$, если его минимальный многочлен над $F$ не имеет кратных корней.

\begin{theorem}[5.1: теорема о примитивном элементе]
Пусть $E=F[\alpha_1,\dots,\alpha_r]$~--- конечное расширение поля $F$, причём $\alpha_2,\dots,\alpha_r$ сепарабельны над
$F$ (а $\alpha_1$~--- не обязательно). Тогда существует $\gamma\in E$ с $E=F[\gamma]$.
\end{theorem}

\begin{proof}
\emph{Конечное $F$.} Тогда $E$ тоже конечно, группа $E^\times$ циклическая, и её образующая $\gamma$ порождает $E$ даже
как кольцо: $E=F[\gamma]$ (\rf{4.19}). Далее считаем $F$ бесконечным.

\emph{Сведение к $r=2$.} Если утверждение доказано для двух образующих, то по индукции
\[
F[\alpha_1,\alpha_2,\alpha_3,\dots,\alpha_r]=F[\alpha_1',\alpha_3,\dots,\alpha_r]=F[\alpha_1'',\alpha_4,\dots,\alpha_r]=\cdots,
\]
где $F[\alpha_1,\alpha_2]=F[\alpha_1']$, $F[\alpha_1',\alpha_3]=F[\alpha_1'']$ и т.\,д. На каждом шаге вторая образующая
($\alpha_2$, $\alpha_3$, \dots) сепарабельна, а про первую ничего не предполагается~--- как раз ситуация случая $r=2$.

\emph{Случай $r=2$.} Пусть $E=F[\alpha,\beta]$, где $\beta$ сепарабелен над $F$. Пусть $f$ и $g$~--- минимальные
многочлены $\alpha$ и $\beta$ над $F$, $s=\deg f$, $t=\deg g$, и $L$~--- поле разложения $fg$ над $E$ (оно же поле
разложения $fg$ над $F$, так как $\alpha,\beta$~--- корни $fg$). Пусть $\alpha_1=\alpha,\alpha_2,\dots,\alpha_{s'}$~---
различные корни $f$ в $L$, а $\beta_1=\beta,\beta_2,\dots,\beta_t$~--- корни $g$; они попарно различны, так как $\beta$
сепарабелен. (С этого места и до конца доказательства индексы у $\alpha_i$ и $\beta_j$ нумеруют корни $f$ и $g$, а не
образующие из формулировки теоремы.)

\emph{Шаг 1: выбор $c$.} При $j\ne1$ имеем $\beta_j\ne\beta$, поэтому уравнение
\[
\alpha_i+X\beta_j=\alpha+X\beta,\qquad\text{т.\,е.}\qquad (\beta-\beta_j)X=\alpha_i-\alpha,
\]
имеет ровно одно решение $X=\dfrac{\alpha_i-\alpha}{\beta-\beta_j}\in L$. Таких <<плохих>> значений конечное число (не
больше $s'(t-1)$), а $F$ бесконечно; выберем $c\in F$, отличное от всех них. Тогда
\[
\alpha_i+c\beta_j\ne\alpha+c\beta\qquad\text{кроме случая } i=1=j
\]
(при $j=1$ равенство $\alpha_i+c\beta=\alpha+c\beta$ означает $\alpha_i=\alpha$, т.\,е.\ $i=1$). Положим
$\gamma=\alpha+c\beta\in E$. Докажем, что $F[\alpha,\beta]=F[\gamma]$; включение $F[\gamma]\subset F[\alpha,\beta]$
очевидно.

\emph{Шаг 2: два многочлена с общим корнем $\beta$.} Многочлены $g(X)\in F[X]$ и $f(\gamma-cX)\in F[\gamma][X]$ имеют
коэффициенты в $F[\gamma]$ и общий корень $\beta$:
\[
g(\beta)=0,\qquad f(\gamma-c\beta)=f(\alpha)=0.
\]
Других общих корней в $L$ у них нет: общий корень~--- это корень $g$, т.\,е.\ некоторое $\beta_j$, причём
$\gamma-c\beta_j$~--- корень $f$, т.\,е.\ $\gamma-c\beta_j=\alpha_i$; по выбору $c$ это возможно лишь при $i=1=j$. Так как
$\beta$~--- простой корень $g$, наибольший общий делитель, вычисленный в $L[X]$, равен
\[
\gcd\bigl(g(X),f(\gamma-cX)\bigr)=X-\beta.
\]

\emph{Шаг 3: $\beta\in F[\gamma]$.} Пусть $h\in F[\gamma][X]$~--- унитарный наибольший общий делитель тех же двух
многочленов, вычисленный в $F[\gamma][X]$. По \rf{2.17} он остаётся наибольшим общим делителем и в $L[X]$, т.\,е.\ $h=X-\beta$.
Коэффициенты $h$ лежат в $F[\gamma]$, поэтому $\beta\in F[\gamma]$, а тогда и $\alpha=\gamma-c\beta\in F[\gamma]$
(здесь важно, что $c\in F$). Итак, $F[\alpha,\beta]=F[\gamma]$.
\end{proof}

Схема доказательства: (1) переходим в поле $L$, где $f$ и $g$ раскладываются; (2) выбираем $c\in F$ так, чтобы все
суммы $\alpha_i+c\beta_j$ с $(i,j)\ne(1,1)$ были отличны от $\gamma=\alpha+c\beta$; (3) тогда у $g(X)$ и $f(\gamma-cX)$ единственный общий корень
$\beta$, наибольший общий делитель $X-\beta$ не зависит от поля, и потому $\beta\in F[\gamma]$.

\begin{remark}[5.2]
Если $F$ бесконечно, то доказательство показывает, что $\gamma$ можно искать в виде
\[
\gamma=\alpha_1+c_2\alpha_2+\dots+c_r\alpha_r,\qquad c_i\in F,
\]
причём на каждом шаге годятся все $c$, кроме конечного числа. Если расширение $E=F[\alpha_1,\dots,\alpha_r]$ ещё и
расширение Галуа, то проверка совсем проста: \emph{элемент $\gamma\in E$ примитивен тогда и только тогда, когда
$\sigma\gamma\ne\gamma$ для каждого $\sigma\ne1$ из $\Gal(E/F)$}. Действительно, $\Gal(E/F[\gamma])$~--- это подгруппа
автоморфизмов, оставляющих $\gamma$ на месте, а по основной теореме $F[\gamma]=E$ ровно тогда, когда эта подгруппа
тривиальна. Именно так мы рассуждали в примере с $\Q[\sqrt2,\sqrt3]$ в начале раздела.

Условия теоремы нельзя ослабить: если несепарабельны хотя бы два из элементов $\alpha_i$, расширение может не быть
простым (пример \rf{5.5}). Чтобы это доказать, нужен ещё один результат.
\end{remark}

\begin{proposition}[5.3]
Пусть $E=F[\gamma]$~--- простое алгебраическое расширение $F$. Тогда промежуточных полей $F\subset M\subset E$ лишь
конечное число.
\end{proposition}

\begin{proof}
\emph{Шаг 1: поле $M$ порождается коэффициентами минимального многочлена $\gamma$ над $M$.} Пусть $F\subset M\subset E$ и
\[
g(X)=X^k+m_{k-1}X^{k-1}+\dots+m_0\in M[X]
\]
~--- минимальный многочлен $\gamma$ над $M$. Положим $M'=F[m_0,\dots,m_{k-1}]$; это поле, так как элементы $m_i\in E$
алгебраичны над $F$. Так как все $m_i\in M$, то $M'\subset M$ и
$g\in M'[X]$. Пусть $g'$~--- минимальный многочлен $\gamma$ над $M'$. Он делит $g$ (в $M'[X]$, а значит, и в $M[X]$),
так как $g(\gamma)=0$. С другой стороны, $g'\in M'[X]\subset M[X]$ и $g'(\gamma)=0$, поэтому $g\mid g'$ в $M[X]$. Оба
многочлена унитарны, значит, $g'=g$. Поскольку $F\subset M'\subset M$, имеем $E=F[\gamma]=M'[\gamma]=M[\gamma]$, а
степень поля $K[\gamma]$ над $K$ равна степени минимального многочлена $\gamma$ над $K$. Следовательно,
\[
[E:M']=\deg g'=\deg g=[E:M],
\]
и по мультипликативности степеней $[M:M']=1$, т.\,е.\ $M=M'=F[m_0,\dots,m_{k-1}]$.
\begin{keepfig}
\begin{center}
\begin{tikzpicture}
\node (A) at (0,0) {$F$};
\node[anchor=base west] (s1) at (A.base east) {$\ \subset\ $};
\node[anchor=base west] (B) at (s1.base east) {$M'=F[m_0,\dots,m_{k-1}]$};
\node[anchor=base west] (s2) at (B.base east) {$\ \subset\ $};
\node[anchor=base west] (C) at (s2.base east) {$M$};
\node[anchor=base west] (s3) at (C.base east) {$\ \subset\ $};
\node[anchor=base west] (D) at (s3.base east) {$E=F[\gamma]$};
\draw[decorate,decoration={brace,amplitude=5pt},cnew] ([yshift=2pt]B.north west) -- node[above=6pt,font=\footnotesize]{$[E:M']=\deg g'=\deg g$} ([yshift=2pt]D.north east);
\draw[decorate,decoration={brace,mirror,amplitude=5pt},cold] ([yshift=-2pt]C.south west) -- node[below=6pt,font=\footnotesize]{$[E:M]=\deg g$} ([yshift=-2pt]D.south east);
\draw[decorate,decoration={brace,mirror,amplitude=5pt},cnew] ([yshift=-34pt]B.south west) -- node[below=6pt,font=\footnotesize]{$[M:M']=1$} ([yshift=-34pt]C.south east);
\end{tikzpicture}
\end{center}
\figcap{Минимальный многочлен $g$ элемента $\gamma$ над $M$ остаётся минимальным и над меньшим полем $M'$, порождённым
его коэффициентами; поэтому $[E:M']=[E:M]$ и $M'=M$.}
\end{keepfig}

\emph{Шаг 2: возможных $g$ конечное число.} Пусть $f$~--- минимальный многочлен $\gamma$ над $F$. Так как $f(\gamma)=0$
и $f\in M[X]$, то $g\mid f$ в $M[X]$, а значит, и в $E[X]$. Пусть $L\supset E$~--- поле разложения $f$; в $L[X]$
многочлен $f=\prod_i(X-\gamma_i)^{e_i}$, и его унитарные делители~--- это произведения $\prod_i(X-\gamma_i)^{d_i}$ с
$0\le d_i\le e_i$; их конечное число. Итак, $g$ пробегает конечное множество унитарных делителей $f$, а $M$ однозначно
определяется своим $g$ (шаг 1). Значит, полей $M$ конечное число.
\end{proof}

Доказательство даёт и описание всех промежуточных полей: каждое из них порождается над $F$ коэффициентами некоторого
унитарного делителя $g$ многочлена $f$ в $E[X]$, имеющего корень $\gamma$. Коэффициенты такого $g$~--- (с точностью до
знака) элементарные симметрические многочлены от части корней $f$, т.\,е.\ <<частично симметрические>> функции от
сопряжённых к $\gamma$: они сохраняются некоторыми, но не обязательно всеми перестановками корней.

\begin{example}[5.3a]
В примере с $\Q[\sqrt2,\sqrt3]$ из начала раздела $\gamma=\sqrt2+\sqrt3$, $f=X^4-10X^2+1$, корни $f$~--- числа $\pm\sqrt2\pm\sqrt3$. Делители $f$ степени $2$
с корнем $\gamma$ получаются, если взять в пару к $\gamma$ один из трёх других корней:
\begin{align*}
(X-\sqrt2-\sqrt3)(X-\sqrt2+\sqrt3)&=X^2-2\sqrt2\,X-1, &&\text{поле } \Q[\sqrt2],\\
(X-\sqrt2-\sqrt3)(X+\sqrt2-\sqrt3)&=X^2-2\sqrt3\,X+1, &&\text{поле } \Q[\sqrt3],\\
(X-\sqrt2-\sqrt3)(X+\sqrt2+\sqrt3)&=X^2-(5+2\sqrt6), &&\text{поле } \Q[\sqrt6].
\end{align*}
Вместе с $f$ (поле $\Q$), с $X-\gamma$ и с тремя делителями степени $3$ (такой делитель равен
\[
f/(X-\gamma')=X^3+\gamma'X^2+(\gamma'^2-10)X+\gamma'(\gamma'^2-10)
\]
для корня $\gamma'\ne\gamma$ многочлена $f$; его коэффициент при $X^2$ равен $\gamma'$, поэтому коэффициенты порождают
$\Q[\gamma']=E$~--- ведь $\gamma'$, корень неприводимого многочлена $f$ степени $4$, имеет степень $4$ над $\Q$) получаем
все промежуточные поля:
$\Q,\ \Q[\sqrt2],\ \Q[\sqrt3],\ \Q[\sqrt6],\ E$~--- в согласии с тем, что у четверной группы ровно пять подгрупп.
\end{example}

Чтобы доказать обращение предложения \rf{5.3}, нужна лемма из линейной алгебры.

\begin{lemma}[5.3b]
Пусть $V$~--- векторное пространство над полем $F$, а $V_1,\dots,V_n$~--- его собственные подпространства
\textup(т.\,е.\ отличные от $V$\textup). Если в $F$ не
меньше $n$ элементов (например, если $F$ бесконечно), то $V\ne V_1\cup\dots\cup V_n$.
\end{lemma}

\begin{proof}
Предположим, что $V=V_1\cup\dots\cup V_n$, и выбросим лишние подпространства: можно считать, что $n$ минимально, т.\,е.\
ни одно $V_i$ не содержится в объединении остальных (иначе его можно выбросить; при этом $n$ только уменьшается, и
условие $|F|\ge n$ сохраняется). Так как $V_1\ne V$, то $n\ge2$. Выберем
\[
u\in V_1,\quad u\notin V_2\cup\dots\cup V_n,\qquad\text{и}\qquad v\in V,\quad v\notin V_1 .
\]
Рассмотрим <<прямую>> $\ell=\{v+cu\mid c\in F\}$, проходящую через $v$ параллельно вектору $u\in V_1$. Так как $u\ne0$ (ведь
$0\in V_2$), отображение $c\mapsto v+cu$ инъективно, и в $\ell$ ровно $|F|$ точек.

\emph{(a) $\ell\cap V_1=\varnothing$.} Если $v+cu\in V_1$, то $v=(v+cu)-cu\in V_1$~--- противоречие.

\emph{(b) $|\ell\cap V_j|\le1$ при $j\ge2$.} Если $v+cu$ и $v+c'u$ лежат в $V_j$ и $c\ne c'$, то
$(c-c')u\in V_j$ и $u\in V_j$~--- противоречие с выбором $u$.

Но $\ell\subset V=V_1\cup\dots\cup V_n$, поэтому по (a) и (b)
\[
\ell=\bigcup_{j=2}^n(\ell\cap V_j),\qquad |F|=|\ell|\le n-1,
\]
что противоречит условию $|F|\ge n$.
\end{proof}

\begin{keepfig}
\begin{center}
\begin{tikzpicture}[scale=1.15,>=stealth]
\fill[cold!10] (-1.6,-0.7) -- (2.6,-0.7) -- (3.5,0.7) -- (-0.7,0.7) -- cycle;
\draw[cold] (-1.6,-0.7) -- (2.6,-0.7) -- (3.5,0.7) -- (-0.7,0.7) -- cycle;
\node[cold] at (3.05,-0.45) {$V_1$};
\coordinate (O) at (0.4,0);
\draw[dashed] (O) -- ($(O)!1.22!(3.7,1.8)$) node[right,font=\footnotesize] {$V_2$};
\draw[dashed] (O) -- ($(O)!1.22!(-0.5,1.8)$) node[left,font=\footnotesize] {$V_3$};
\fill (O) circle (1.2pt) node[below=1pt] {$0$};
\draw[->,thick,cold] (O) -- (1.5,0) node[below] {$u$};
\draw[->,thick] (O) -- (0.9,1.8);
\draw[cnew,line width=1.1pt] (-1.3,1.8) -- (4.6,1.8) node[right] {$\ell=v+Fu$};
\fill[cnew] (0.9,1.8) circle (1.5pt) node[above,font=\footnotesize] {$v$};
\foreach \x/\l in {2.0/v+u,3.1/v+2u}{\fill[cnew] (\x,1.8) circle (1.5pt); \node[cnew,above,font=\footnotesize] at (\x,1.83) {$\l$};}
\fill (3.7,1.8) circle (1.8pt);
\fill (-0.5,1.8) circle (1.8pt);
\end{tikzpicture}
\end{center}
\figcap{Доказательство леммы \rf{5.3b}: прямая $\ell=v+Fu$ (красная) параллельна плоскости $V_1$ и не пересекает её, а с
каждым другим подпространством $V_j$ (пунктир) имеет не более одной общей точки (чёрные точки). Поэтому $\ell$, а с ней и
$V$, не покрываются подпространствами $V_1,\dots,V_n$, если в $\ell$ больше $n-1$ точек.}
\end{keepfig}

Условие $|F|\ge n$ ослабить нельзя: плоскость $\F_q^2$ есть объединение всех прямых, проходящих через $0$ (каждый
вектор $v$ лежит на прямой $\F_qv$), а таких прямых $(q^2-1)/(q-1)=q+1$ (каждая содержит $q-1$ ненулевых векторов, и
различные прямые пересекаются лишь в $0$). Значит, при $n=q+1=|F|+1$ утверждение леммы неверно.

\begin{corollary}[5.3c]
Пусть $F$ бесконечно, $L\subset V$~--- подпространство и $L\subset K_1\cup\dots\cup K_n$ для подпространств $K_i\subset V$.
Тогда $L\subset K_i$ для некоторого $i$.
\end{corollary}

\begin{proof}
Иначе $L\cap K_1,\dots,L\cap K_n$~--- собственные подпространства $L$, объединение которых равно $L$, что невозможно по
лемме \rf{5.3b}.
\end{proof}

\begin{remark}
Для бесконечного поля $F$ и конечномерного $V$ лемму можно доказать и иначе. Отождествим $V$ с $F^m$ и выберем для каждого
$i$ ненулевую линейную форму $f_i$, обращающуюся в нуль на $V_i$. Если $V=\bigcup V_i$, то многочлен $f_1\cdots f_n$
обращается в нуль во всех точках $F^m$, а значит, равен нулю как многочлен (\rf{5.19})~--- но произведение ненулевых
многочленов не равно нулю.
\end{remark}

\begin{remark}[5.4]
У предложения \rf{5.3} есть обращение: \emph{если $E/F$ конечно и промежуточных полей $F\subset M\subset E$ конечное число, то
$E/F$~--- простое расширение.}

\emph{Доказательство.} Если $F$ конечно, то конечно и $E$, и $E/F$ просто по \rf{4.19}. Пусть $F$ бесконечно, а
$M_1,\dots,M_k$~--- все промежуточные поля, отличные от $E$. Это собственные $F$-подпространства $E$, и по лемме \rf{5.3b}
найдётся $\gamma\in E$, не лежащий ни в одном из них. Поле $F[\gamma]$ промежуточное и не совпадает ни с одним $M_i$
(так как $\gamma\notin M_i$); значит, $F[\gamma]=E$.

Это даёт второе доказательство теоремы \rf{5.1} для сепарабельного $E/F$: такое расширение содержится в конечном
расширении Галуа $\tilde E/F$~--- например, в поле разложения над $E$ произведения $h$ \emph{различных} минимальных
многочленов образующих. Это поле содержит $E$ и является полем разложения $h$ над $F$ (образующие~--- корни $h$), а $h$
сепарабелен: его множители сепарабельны, а различные унитарные неприводимые многочлены не имеют общих корней (\rf{2.17});
поэтому $\tilde E/F$~--- расширение Галуа (\rf{3.10}). Промежуточные поля $\tilde E/F$ по основной теореме соответствуют
подгруппам конечной группы $\Gal(\tilde E/F)$, и их конечное число. Тем более конечно число промежуточных полей $E/F$, и $E/F$ просто.
\end{remark}

\begin{example}[5.5]
Простейшее непростое конечное расширение. Пусть $k$~--- поле характеристики $p$, $E=k(X,Y)$~--- поле рациональных
функций от двух переменных и $F=k(X^p,Y^p)$. Покажем, что $[E:F]=p^2$, но $E\ne F[\alpha]$ ни для какого $\alpha$
(последнее~--- двумя способами).

\emph{Степень $[E:F]=p^2$.} Рассмотрим башню $F\subset F[X]\subset F[X][Y]=E$.
\begin{itemize}[leftmargin=*]
\item $[F[X]:F]=p$. Элемент $X$~--- корень многочлена $T^p-X^p\in F[T]$, а он неприводим над $F$ по \rf{2.3}, так как
$X^p$ не является $p$-й степенью в $F$: если бы $X^p=\beta^p$ с $\beta\in F$, то $X=\beta\in F$ (извлечение
корня $p$-й степени в характеристике $p$ однозначно: из $x^p=y^p$ следует $(x-y)^p=0$). А $X\notin F$: элемент $F$ имеет вид $a(X^p,Y^p)/b(X^p,Y^p)$ с
многочленами $a,b$, $b\ne0$, а в равенстве $X\,b(X^p,Y^p)=a(X^p,Y^p)$ все мономы слева имеют показатель при $X$,
сравнимый с $1$ по модулю $p$, а все мономы справа~--- делящийся на $p$; так как левая часть ненулевая, это невозможно.
\item $[E:F[X]]=p$. Поле $F[X]=k(X^p,Y^p)(X)$ совпадает с $k(X,Y^p)$ (оба~--- наименьшее подполе $E$, содержащее $k$,
$X$ и $Y^p$). Элемент $Y$~--- корень $T^p-Y^p$, и тем же рассуждением о показателях (теперь при $Y$) $Y\notin k(X,Y^p)$;
значит, $T^p-Y^p$ неприводим над $F[X]$ по \rf{2.3}.
\end{itemize}
По мультипликативности степеней $[E:F]=p\cdot p=p^2$.

\emph{Первый способ: элементы малой степени.} Пусть $\alpha\in E$, $\alpha=a(X,Y)/b(X,Y)$. В характеристике $p$
возведение в степень $p$~--- гомоморфизм, поэтому $\alpha^p=a^{(p)}(X^p,Y^p)/b^{(p)}(X^p,Y^p)$, где у $a^{(p)},b^{(p)}$
коэффициенты~--- $p$-е степени коэффициентов $a,b$. Значит, $\alpha^p\in F$, и $\alpha$~--- корень многочлена
$T^p-\alpha^p\in F[T]$. Минимальный многочлен $\alpha$ делит его, так что
\[
[F[\alpha]:F]\le p<p^2=[E:F],
\]
и $F[\alpha]\ne E$. Примитивного элемента нет.

\emph{Второй способ: бесконечно много промежуточных полей.} Для $c\in F$ положим $M_c=F[X+cY]$. Тогда:
\begin{itemize}[leftmargin=*]
\item $[M_c:F]\le p$, так как $(X+cY)^p=X^p+c^pY^p\in F$;
\item $[E:M_c]\le p$, так как $E=M_c[Y]$ (ведь $X=(X+cY)-cY$ и $c\in F\subset M_c$), а $Y^p\in F\subset M_c$;
\item $[M_c:F]\cdot[E:M_c]=[E:F]=p^2$.
\end{itemize}
Отсюда $[M_c:F]=[E:M_c]=p$: степени в башне $F\subset M_c\subset E$ равны $p$, $p$ и $p^2$. Далее, при $c\ne c'$ поля
$M_c$ и $M_{c'}$ различны: если $M_c=M_{c'}$, то это поле содержит $X+cY$ и $X+c'Y$, а с ними $(c-c')Y$, т.\,е.\ $Y$
(так как $c-c'\in F^\times\subset M_c$), и $X$; тогда $M_c=F[X,Y]=E$, что противоречит $[M_c:F]=p<p^2$. Поле $F$
бесконечно (в нём лежат различные элементы $X^p,X^{2p},X^{3p},\dots$), поэтому промежуточных полей бесконечно много при
любом $k$, даже при $k=\F_p$, и по предложению \rf{5.3} расширение $E/F$ не простое.
\begin{keepfig}
\begin{center}
\begin{tikzpicture}[yscale=1.2]
\node (E) at (-0.6,2.4) {$E=k(X,Y)$};
\node (F) at (-0.6,0) {$F=k(X^p,Y^p)$};
\node[cnew] (M1) at (-5.0,1.2) {$F[X]=M_0$};
\node[cnew] (M2) at (-2.3,1.2) {$M_c=F[X+cY]$};
\node[cnew] (M3) at (0.1,1.2) {$M_{c'}$};
\node[cnew] (M4) at (1.4,1.2) {$\cdots$};
\node[cnew] (M5) at (2.7,1.2) {$F[Y]$};
\foreach \k in {1,3,5}{\draw (E) -- (M\k); \draw (F) -- (M\k);}
\draw (E) -- node[pos=0.5,below right,font=\scriptsize]{$p$} (M2);
\draw (F) -- node[pos=0.5,above right,font=\scriptsize]{$p$} (M2);
\draw[decorate,decoration={brace,amplitude=5pt},cold] (4.1,2.4) -- node[right=6pt,font=\footnotesize]{$p^2$} (4.1,0);
\end{tikzpicture}
\end{center}
\figcap{Расширение $k(X,Y)/k(X^p,Y^p)$ степени $p^2$: каждое из полей $M_c=F[X+cY]$ ($c\in F$) имеет степень $p$ над
$F$, и $E$ имеет над ним степень $p$; разным $c$ отвечают разные поля. Поле $F[Y]$ тоже промежуточное, но не имеет вида
$M_c$ (иначе $X=(X+cY)-cY\in F[Y]$ и $F[Y]=E$); $F[X]=M_0$.}
\end{keepfig}
\end{example}
